\begin{tikzpicture}
\begin{axis}[
    width=0.8\textwidth,
    ybar, 
    ymin=0,
    % Define exactly where the ticks go
    xtick={0, 1}, 
    xticklabels={AdS ($V<0$), dS ($V>0$)},
    % Styling to make it look professional
    ylabel={Vacíos encontrados},
    xlabel={Tipo de vacío},
    title={Actually Lifted Vacua},
    enlarge x limits=0.6,
    nodes near coords, % This will print the number of vacua on top of each bar
]
\addplot [
    hist={
        bins=2,
        data min=-1.5, % AdS (0) will be in the middle of [-0.5, 0.5]
        data max=1.5   % dS (1) will be in the middle of [0.5, 1.5]
    }
] table [
    y expr={\thisrow{V} < 0 ? 0 : 1}
] {\losDatos};
\end{axis}
\end{tikzpicture}

\begin{tikzpicture}
\small
\begin{axis}[
    ybar interval, % Classic histogram look
    % Si 0, etiqueta estable; si 1, etiqueta taquiónico
    xticklabels={AdS ($V<0$), dS ($V>0$)},
    xticklabel style={rotate=45, anchor=north east, font=\tiny},
    xlabel={Cantidad de vacíos encontrados por tipo después del levantamiento},
    ylabel={Vacíos encontrados},
    title={Distribution of Found Vacua},
    scale=0.7
]
\addplot [
    fill=blue!30, 
    draw=black,
    hist={
        bins=2, % Adjust this to see more/less detail
        data min=-1, % Optional: focus on a specific range
        data max= 1
    }
] table [
    y expr={\thisrow{V} < 0 ? 0 : 1}
] {\losDatos};
\end{axis}
\end{tikzpicture}
